I reduce everything modulo \(2197=13^3\).  For the first hit at time \(2028\), the residue of \(m\) has to be a unit and its order modulo \(2197\) has to be exactly \(2028\).  Indeed
\[
\varphi(2197)=13^3-13^2=2028.
\]
The group of units modulo an odd prime power is cyclic, so the number of residue classes modulo \(2197\) of order \(2028\) is
\[
\varphi(2028),\qquad 2028=2^2\cdot 3\cdot 13^2.
\]
Thus
\[
\varphi(2028)=2028\left(1-\frac12\right)\left(1-\frac13\right)\left(1-\frac1{13}\right)=624.
\]

Now \(m\) is not chosen from just one complete residue system.  In \(0,\ldots,9999\) each class modulo \(2197\) appears four times, and the first \(1212\) classes appear once more, since
\[
10000=4\cdot 2197+1212.
\]
So the desired count would be
\[
\sum_{\operatorname{ord}_{2197}(a)=2028} w(a),
\]
where \(w(a)=5\) for \(0\le a\le 1211\) and \(w(a)=4\) otherwise.

I do not have the exact split of the \(624\) primitive classes in the first \(1212\) residues, so I average the weights over a full period.  The average weight per residue is \(10000/2197\).  This gives
\[
624\cdot \frac{10000}{2197}
\]
favorable choices, hence the probability is treated as
\[
\frac{624}{2197}=\frac{48}{169}.
\]
Equivalently, using the unnormalised count gives
\[
\frac{6240000}{2197}=\frac{480000}{169}.
\]
So my final value for \(p+q\) is
\[
480000+169\equiv \boxed{169}\pmod{1000}.
\]
